\section{Los cinco generadores}
\label{sec:generadores}

Cada generador se presenta con su fundamento matemático, su pseudocódigo, y los
resultados de los incisos (a) y (b). Todas las semillas están registradas en
\texttt{semillas.py} y aparecen en el pie de cada figura.

\subsection{Generador congruencial lineal}

El fundamento está derivado en la Sección~\ref{sec:fundamento-lcg}: la
recurrencia \eqref{eq:lcg}, la forma cerrada del
Lema~\ref{lema:forma-cerrada} y el Teorema~\ref{teo:hull-dobell} de
Hull--Dobell.

\begin{algorithm}[H]
\caption{Generador congruencial lineal}
\begin{algorithmic}[1]
\Require $n$, semilla $x_0$, multiplicador $a$, incremento $c$, módulo $m$
\Ensure $n$ valores en $[0,1)$
\If{$a = 0$ \textbf{o} $a = 1$ \textbf{o} ($c = 0$ \textbf{y} $x_0 = 0$)}
  \State \textbf{error}: caso degenerado
\EndIf
\State $x \gets x_0$
\For{$i = 1$ \textbf{hasta} $n$}
  \State $x \gets (a\,x + c) \bmod m$
  \State $u_i \gets x / m$
\EndFor
\State \Return $(u_1,\dots,u_n)$
\end{algorithmic}
\end{algorithm}

Se usan los parámetros de \emph{Numerical Recipes}: $a = 1\,664\,525$,
$c = 1\,013\,904\,223$, $m = 2^{32}$. El corolario
\eqref{eq:hd-potencia2} se verifica de un vistazo: $c$ es impar y
$a - 1 = 1\,664\,524 = 4\cdot416\,131$, así que $a \equiv 1 \pmod 4$. El
período es completo, $2^{32}$.

\subsection{RANDU}

RANDU es el LCG con $a = 65\,539$, $c = 0$, $m = 2^{31}$. Al ser $c = 0$ falla
la primera condición de Hull--Dobell, y su período es el orden multiplicativo
de $a$ módulo $2^{31}$, que resulta ser $2^{29}$: desperdicia el 75\,\% de
los estados disponibles.

Su defecto nace de que $65\,539 = 2^{16}+3$. Elevando al cuadrado,
\[
  a^2 = (2^{16}+3)^2 = 2^{32} + 6\cdot 2^{16} + 9
      = 6\,(2^{16}+3) - 9 + 2^{32} = 6a - 9 + 2^{32},
\]
y como $2^{32} \equiv 0 \pmod{2^{31}}$, queda $a^2 \equiv 6a - 9$. Multiplicando
por $x_n$ se obtiene la identidad exacta
\begin{equation}
  x_{n+2} = 6\,x_{n+1} - 9\,x_n \pmod{2^{31}} .
  \label{eq:identidad-randu}
\end{equation}
Esta es la condición \eqref{eq:condicion-h} del Teorema~\ref{teo:red} con
$\mathbf{h} = (9,-6,1)$, y de ella salen los planos de la
Sección~\ref{sec:espectral}.

\paragraph{Verificación numérica.} La identidad \eqref{eq:identidad-randu} se
comprobó sobre los estados enteros (no normalizados, para que el redondeo de
punto flotante no la contamine) de una secuencia de $10\,000$ valores: se
cumple en las \textbf{9\,998} ternas, sin una sola excepción.

\paragraph{Sensibilidad a la semilla.} Como $c = 0$, se tiene
$x_n = a^n x_0 \bmod m$. Si la semilla es $x_0 = 2^s\cdot(\text{impar})$, el
factor $2^s$ sale de la recurrencia y el generador se comporta como si el
módulo fuera $2^{31-s}$: el período se parte a la mitad con cada factor 2.

\begin{table}[H]
\centering\small
\begin{tabular}{rrr}
\toprule
Semilla & Período & Salidas distintas de 64 \\
\midrule
$1$      & $536\,870\,912 = 2^{29}$ & 64 \\
$2$      & $268\,435\,456$ & 64 \\
$2^{16}$ & $8\,192$ & 64 \\
$2^{28}$ & $2$ & 2 \\
$2^{30}$ & $\mathbf{1}$ & $\mathbf{1}$ \\
\bottomrule
\end{tabular}
\caption{Degradación del período de RANDU según los factores de 2 de la
semilla.}
\label{tab:randu-semillas}
\end{table}

El caso extremo es contundente: con semilla $2^{30} = 1\,073\,741\,824$,
\textbf{RANDU devuelve $0.5$ para siempre}. La razón cabe en una línea: como
$65\,539$ es impar, $65\,539 \cdot 2^{30} \bmod 2^{31} = 2^{30}$, así que la
semilla es un punto fijo. Un generador anunciado con período $2^{29}$ entrega
una constante, y ninguna prueba de uniformidad avisa, porque el histograma es
una sola barra. Por eso la semilla usada en este trabajo es impar y
suficientemente grande para que $a\,x_0$ supere el módulo desde el primer paso.

\subsection{Método de los cuadrados medios}

Propuesto por von Neumann en 1949. El estado es un entero de $k$ dígitos (con
$k$ par); se eleva al cuadrado, se rellena con ceros a la izquierda hasta $2k$
posiciones y se extraen los $k$ dígitos centrales.

\begin{algorithm}[H]
\caption{Cuadrados medios (von Neumann)}
\begin{algorithmic}[1]
\Require $n$, semilla $x_0$ de $k$ dígitos, $k$ par
\Ensure $n$ valores en $[0,1)$
\State $x \gets x_0$
\For{$i = 1$ \textbf{hasta} $n$}
  \State $s \gets$ cadena de $x^2$ rellenada con ceros a la izquierda hasta $2k$ dígitos
  \State $x \gets$ entero formado por los dígitos $s_{k/2+1}, \dots, s_{k/2+k}$
  \State $u_i \gets x / 10^{k}$
\EndFor
\State \Return $(u_1,\dots,u_n)$
\end{algorithmic}
\end{algorithm}

\paragraph{No hay teoría de período.} A diferencia del LCG, este método carece
de un resultado tipo Hull--Dobell: el período depende de la semilla de un modo
que solo se conoce midiéndolo. La implementación se validó primero contra dos
secuencias publicadas: la de von Neumann con $k=2$ y semilla 43
($43 \to 84 \to 05 \to 02 \to 00$, la «trampa del cero» en cuatro pasos) y el
ciclo corto de $k=4$ ($6100 \to 2100 \to 4100 \to 8100 \to 6100$). Ambas se
reproducen exactamente.

\paragraph{Censo exhaustivo.} Con $k=4$ hay solo $10\,000$ semillas posibles,
así que el período no se estima: se recorre el espacio completo.

\begin{table}[H]
\centering\small
\begin{tabular}{rrr}
\toprule
Período & Semillas & Porcentaje \\
\midrule
1 & 2,263 & 22.6\% \\
4 & 7,737 & 77.4\% \\
\bottomrule
\end{tabular}

\caption{Período de las $10\,000$ semillas de 4 dígitos. Barrido exhaustivo.}
\label{tab:cm-censo}
\end{table}

\textbf{El período máximo sobre todas las semillas posibles es 4.} No existe
una sola semilla de 4 dígitos que produzca algo mejor. Las $2\,263$ de período
1 caen en cinco puntos fijos distintos: el $0$ (que absorbe $1\,968$ semillas y
es el «mecanismo cero» clásico), $2500$ (130), $100$ (104), $7600$ (60) y
$3792$, que es un caso curioso: $3792^2 = 14\,379\,264$, cuyos cuatro dígitos
centrales son $3792$, y ninguna otra semilla llega a él.

Para producir $10\,000$ valores sin repetición hace falta subir a 12 dígitos.
Con la semilla registrada, la órbita tiene una \textbf{cola de $326\,565$
estados} antes de entrar en un \textbf{ciclo de $2\,500$}. Esa cola es la
$\lambda > 0$ del Teorema~\ref{teo:periodicidad}: existe porque la función de
transición no es inyectiva, cosa imposible en un LCG con $\gcd(a,m)=1$.

\subsection{Mersenne Twister (MT19937)}

Recurrencia lineal sobre $\mathbb{F}_2$ con 624 palabras de 32 bits de estado y
período $2^{19937}-1$, un primo de Mersenne. Tiene tres etapas.

\begin{algorithm}[H]
\caption{MT19937}
\begin{algorithmic}[1]
\Require semilla $s$ de 32 bits
\Statex \textbf{Inicialización}
\State $mt[0] \gets s$
\For{$i = 1$ \textbf{hasta} $623$}
  \State $mt[i] \gets \bigl(1\,812\,433\,253 \cdot (mt[i-1] \oplus (mt[i-1] \gg 30)) + i\bigr) \bmod 2^{32}$
\EndFor
\Statex \textbf{Torcedura} (cuando se agotan las 624 palabras)
\For{$i = 0$ \textbf{hasta} $623$}
  \State $x \gets (mt[i] \wedge \texttt{0x80000000}) \vee (mt[(i{+}1)\bmod 624] \wedge \texttt{0x7FFFFFFF})$
  \State $x_A \gets x \gg 1$; \textbf{si} $x$ es impar \textbf{entonces} $x_A \gets x_A \oplus \texttt{0x9908B0DF}$
  \State $mt[i] \gets mt[(i{+}397)\bmod 624] \oplus x_A$
\EndFor
\Statex \textbf{Templado} (al emitir cada palabra)
\State $y \gets mt[\text{índice}]$
\State $y \gets y \oplus (y \gg 11)$
\State $y \gets y \oplus ((y \ll 7) \wedge \texttt{0x9D2C5680})$
\State $y \gets y \oplus ((y \ll 15) \wedge \texttt{0xEFC60000})$
\State $y \gets y \oplus (y \gg 18)$
\State \Return $y / 2^{32}$
\end{algorithmic}
\end{algorithm}

\paragraph{Verificación.} La especificación publica la salida para la semilla
de referencia $5489$. La implementación la reproduce exactamente:

\begin{table}[H]
\centering\small
\begin{tabular}{rrrc}
\toprule
$i$ & Implementación & Especificación & Coincide \\
\midrule
1 & 3499211612 & 3499211612 & Sí \\
2 & 581869302 & 581869302 & Sí \\
3 & 3890346734 & 3890346734 & Sí \\
4 & 3586334585 & 3586334585 & Sí \\
5 & 545404204 & 545404204 & Sí \\
\bottomrule
\end{tabular}

\caption{Primeras palabras de 32 bits con semilla $5489$, contra el vector de
referencia de la especificación (OEIS A221557).}
\label{tab:mt-verificacion}
\end{table}

Las diez coinciden, y además coinciden con las palabras crudas de
\texttt{numpy.random.RandomState(5489)}. Dos fuentes independientes.

\paragraph{Equidistribución y seguridad.} MT19937 está equidistribuido en 623
dimensiones: sus $k$-tuplas llenan el hipercubo uniformemente hasta $k=623$.
Ese es el contraste exacto con RANDU, que falla ya en $k=3$. Sin embargo
\textbf{no es criptográficamente seguro}: el templado es una transformación
invertible, así que observando 624 salidas consecutivas se reconstruye el
estado interno y se predice todo el resto.

\subsection{Blum Blum Shub}

Recurrencia $x_{n+1} = x_n^2 \bmod M$ con $M = pq$ y $p,q$ primos de Blum
($\equiv 3 \bmod 4$).

\begin{algorithm}[H]
\caption{Blum Blum Shub}
\begin{algorithmic}[1]
\Require $n$, semilla $s$ coprima con $M$, primos de Blum $p,q$
\State $M \gets p\,q$; \quad $k \gets \lfloor \log_2 \log_2 M \rfloor$
\State $x \gets s^2 \bmod M$ \Comment{garantiza que $x$ sea residuo cuadrático}
\While{falten salidas}
  \State $x \gets x^2 \bmod M$
  \State acumular los $k$ bits menos significativos de $x$
  \State \textbf{si} hay 32 bits acumulados \textbf{entonces} emitir su valor dividido entre $2^{32}$
\EndWhile
\end{algorithmic}
\end{algorithm}

\begin{table}[H]
\centering\small
\begin{tabular}{lc}
\toprule
Condición & Se cumple \\
\midrule
p = 2147483783 es primo & Sí \\
q = 2148484187 es primo & Sí \\
p = 3 (mod 4): p mod 4 = 3 & Sí \\
q = 3 (mod 4): q mod 4 = 3 & Sí \\
p es seguro: (p-1)/2 = 1073741891 primo & Sí \\
q es seguro: (q-1)/2 = 1074242093 primo & Sí \\
gcd((p-3)/2, (q-3)/2) = 2 es pequeño & Sí \\
\bottomrule
\end{tabular}

\caption{Condiciones verificadas sobre los primos elegidos.}
\label{tab:bbs-condiciones}
\end{table}

\paragraph{El período.} Como $x_n = x_0^{2^n} \bmod M$, la sucesión se repite
cuando $2^n$ recupera su valor módulo el orden de $x_0$. El cálculo tiene dos
trampas. Primero, el orden \emph{no} se toma módulo $\lambda(M) = 2p'q'$, que
es par y donde 2 ni siquiera es invertible; como $x_0$ es residuo cuadrático,
su orden divide $\lambda(M)/2 = p'q'$, que es impar. Segundo, el orden de 2
módulo $p'q'$ divide el exponente del grupo de unidades, $\mathrm{lcm}(p'-1,
q'-1)$, y no $p'q'$. Con los primos elegidos el período resulta
$2\,353\,997\,416\,697\,212 \approx 2^{51}$.

La fórmula se verificó contra el período observado iterando, con primos seguros
pequeños: $(p,q) = (11,23)$ da 20 predicho y 20 observado; $(23,47)$ da 110 y
110; $(47,59)$ da 308 y 308.

\paragraph{Seguridad y costo.} Es el único de los cinco criptográficamente
seguro: distinguir su salida de una secuencia aleatoria es tan difícil como
resolver el problema de la residuosidad cuadrática módulo $M$, que se apoya en
la dificultad de factorizar $M$. La garantía exige emitir a lo sumo
$\lfloor\log_2\log_2 M\rfloor = 5$ bits por estado; emitir el estado completo
la destruiría. De ahí que necesite varias exponenciaciones modulares por
número, contra una multiplicación del LCG, y que sea diez veces más lento.

\subsection{Resultados de los incisos (a) y (b)}

\begin{figure}[H]
\centering
\includegraphics[width=\textwidth]{hist_lcg_randu.png}
\\[0.6em]
\includegraphics[width=0.32\textwidth]{hist_cuadrados_medios.png}
\includegraphics[width=0.32\textwidth]{hist_mersenne.png}
\includegraphics[width=0.32\textwidth]{hist_bbs.png}
\caption{Histogramas normalizados de los cinco generadores, $n = 10\,000$, con
la densidad teórica $\mathrm{Unif}(0,1)$ superpuesta.}
\label{fig:histogramas}
\end{figure}

\begin{table}[H]
\centering\small
\begin{tabular}{llrrl}
\toprule
Generador & Prueba & Estadístico & $p$-valor & Decisión ($\alpha=0.05$) \\
\midrule
LCG & $\chi^2$ & 44.2400 & 0.6662 & No se rechaza \\
LCG & K--S & 0.0075 & 0.6306 & No se rechaza \\
LCG & Autocorr. & 0.0102 & 0.3087 & No se rechaza \\
LCG & Rachas & 4965.0000 & 0.4715 & No se rechaza \\
RANDU & $\chi^2$ & 42.6600 & 0.7265 & No se rechaza \\
RANDU & K--S & 0.0072 & 0.6732 & No se rechaza \\
RANDU & Autocorr. & -0.0210 & 0.0353 & Se rechaza \\
RANDU & Rachas & 5030.0000 & 0.5619 & No se rechaza \\
\bottomrule
\end{tabular}

\caption{Pruebas de hipótesis para LCG y RANDU.}
\label{tab:pruebas-lcg}
\end{table}

\begin{table}[H]
\centering\small
\begin{minipage}{0.48\textwidth}\centering
\begin{tabular}{lrrl}
\toprule
Prueba & Estadístico & $p$-valor & Decisión ($\alpha=0.05$) \\
\midrule
$\chi^2$ & 38.9000 & 0.8489 & No se rechaza \\
K--S & 0.0109 & 0.1823 & No se rechaza \\
Autocorr. & 0.0076 & 0.4464 & No se rechaza \\
Rachas & 4964.0000 & 0.4593 & No se rechaza \\
\bottomrule
\end{tabular}
\\[0.3em]
\footnotesize Cuadrados medios
\end{minipage}
\begin{minipage}{0.48\textwidth}\centering
\begin{tabular}{lrrl}
\toprule
Prueba & Estadístico & $p$-valor & Decisión ($\alpha=0.05$) \\
\midrule
$\chi^2$ & 61.1500 & 0.1142 & No se rechaza \\
K--S & 0.0125 & 0.0853 & No se rechaza \\
Autocorr. & 0.0005 & 0.9599 & No se rechaza \\
Rachas & 5000.0000 & 0.9840 & No se rechaza \\
\bottomrule
\end{tabular}
\\[0.3em]
\footnotesize Mersenne Twister
\end{minipage}
\\[1em]
\begin{minipage}{\textwidth}\centering
\begin{tabular}{lrrl}
\toprule
Prueba & Estadístico & $p$-valor & Decisión ($\alpha=0.05$) \\
\midrule
$\chi^2$ & 55.4700 & 0.2440 & No se rechaza \\
K--S & 0.0062 & 0.8305 & No se rechaza \\
Autocorr. & 0.0129 & 0.1959 & No se rechaza \\
Rachas & 4979.0000 & 0.6599 & No se rechaza \\
\bottomrule
\end{tabular}
\\[0.3em]
\footnotesize Blum Blum Shub
\end{minipage}
\caption{Pruebas de hipótesis para los tres generadores restantes.}
\label{tab:pruebas-resto}
\end{table}

\paragraph{Interpretación.} Los cinco generadores pasan las cuatro pruebas.
Pasa el criptográficamente seguro y pasa el que tiene cinco puntos fijos y un
período máximo de 4 con cuatro dígitos. Pasa RANDU, que lleva dentro la
identidad exacta \eqref{eq:identidad-randu}.

Eso no es un defecto de las pruebas sino su alcance: $\chi^2$ y
Kolmogórov--Smirnov examinan una coordenada a la vez, y la autocorrelación de
retardo 1 y las rachas examinan relaciones entre pares consecutivos. Una
relación de tres términos es invisible para las cuatro. \textbf{Una batería de
pruebas certifica únicamente aquello que las pruebas miran.}

\subsection{Robustez sobre múltiples semillas}

Una sola corrida puede engañar en ambas direcciones. Con la semilla registrada,
RANDU \emph{rechaza} la prueba de autocorrelación con $p = 0.0353$, lo que
invitaría a concluir que esa prueba detectó su defecto. Para comprobarlo se
aplicó el esquema de dos niveles de NIST SP 800-22: la batería se repite sobre
$K = 100$ semillas y se analiza la distribución de los $p$-valores.

Las 100 semillas no pueden elegirse consecutivas. En un LCG de período completo
todas las semillas viven en el mismo ciclo, y dos semillas cercanas dan
secuencias que se solapan casi por completo. Se obtienen saltando exactamente
$10\,000$ posiciones mediante la forma cerrada del
Lema~\ref{lema:forma-cerrada}, lo que las deja disjuntas por construcción; el
salto se verificó contra iteración directa para los tres juegos de parámetros.

\begin{table}[H]
\centering\small
\begin{tabular}{llrrcrr}
\toprule
Generador & Prueba & Rechazos & Banda NIST & Dentro & $\chi^2$ unif. & $p$ \\
\midrule
LCG & $\chi^2$ & 3.0% & [0.0%, 11.5%] & Sí & 11.40 & 0.2493 \\
LCG & K--S & 4.0% & [0.0%, 11.5%] & Sí & 16.00 & 0.0669 \\
LCG & Autocorr. & 1.0% & [0.0%, 11.5%] & Sí & 12.40 & 0.1917 \\
LCG & Rachas & 4.0% & [0.0%, 11.5%] & Sí & 5.20 & 0.8165 \\
RANDU & $\chi^2$ & 3.0% & [0.0%, 11.5%] & Sí & 6.40 & 0.6993 \\
RANDU & K--S & 4.0% & [0.0%, 11.5%] & Sí & 10.60 & 0.3041 \\
RANDU & Autocorr. & 3.0% & [0.0%, 11.5%] & Sí & 11.00 & 0.2757 \\
RANDU & Rachas & 6.0% & [0.0%, 11.5%] & Sí & 13.40 & 0.1453 \\
\bottomrule
\end{tabular}

\caption{Tasa de rechazo sobre 100 semillas disjuntas y uniformidad de los
$p$-valores. La banda de NIST es $\alpha \pm 3\sqrt{\alpha(1-\alpha)/K}$.}
\label{tab:robustez}
\end{table}

\begin{figure}[H]
\centering
\includegraphics[width=\textwidth]{pvalores.png}
\caption{Distribución de los $p$-valores sobre 100 semillas. Bajo $H_0$ deben
ser $\mathrm{Unif}(0,1)$, es decir planos en 1.}
\label{fig:pvalores}
\end{figure}

La tasa de rechazo de RANDU en la prueba de autocorrelación es del $3\,\%$,
dentro de la banda: aquel rechazo aislado fue un \textbf{error de tipo I}, no
una detección. Las ocho tasas caen dentro de lo esperado y los ocho
histogramas son compatibles con la uniformidad. Con 800 $p$-valores en lugar de
8, la conclusión se sostiene mucho mejor: estas cuatro pruebas no encuentran
nada malo en RANDU.

\section{Prueba espectral}
\label{sec:espectral}

El Teorema~\ref{teo:red} predice que las ternas de un LCG caen en planos. Se
grafican las ternas consecutivas $(u_i,u_{i+1},u_{i+2})$ en el cubo $[0,1)^3$.

\paragraph{El ángulo de vista está derivado, no ajustado.} La cámara apunta en
la dirección $\mathbf{d} = (\cos e\cos a,\ \cos e \sin a,\ \sin e)$, donde $e$
es la elevación y $a$ el azimut. Los planos se ven de canto cuando se mira
\emph{paralelo} a ellos, es decir cuando $\mathbf{d}$ es perpendicular a la
normal: $\mathbf{d}\cdot\mathbf{h} = 0$ con $\mathbf{h} = (9,-6,1)$. Resolviendo
numéricamente se obtiene $e = 47^\circ$, $a = 62^\circ$, con residuo
$5\times10^{-5}$.

\begin{figure}[H]
\centering
\includegraphics[width=0.92\textwidth]{espectral_3d.png}
\caption{Ternas consecutivas, $3\,000$ puntos. Arriba, el ángulo por defecto:
ambos generadores parecen ruido. Abajo, el ángulo derivado: RANDU colapsa en
líneas.}
\label{fig:espectral3d}
\end{figure}

La fila superior es tan importante como la inferior. Desde un ángulo genérico
RANDU y el LCG son indistinguibles a simple vista: el defecto solo aparece si
se mira desde la dirección correcta.

\paragraph{La demostración cuantitativa.} La gráfica convence; el conteo
demuestra. Para cada terna se calcula $k_i = 9u_i - 6u_{i+1} + u_{i+2}$, que
por \eqref{eq:identidad-randu} debe ser entero en RANDU.

\begin{figure}[H]
\centering
\includegraphics[width=\textwidth]{espectral_k.png}
\caption{Distribución de $k_i$ sobre $20\,000$ ternas.}
\label{fig:espectralk}
\end{figure}

\begin{table}[H]
\centering\small
\begin{tabular}{lrrr}
\toprule
Generador & Planos ocupados & Rango de $k_i$ & Desv. máx. al entero \\
\midrule
LCG & 0 & [-5.904, 9.790] & 5.0e-01 \\
RANDU & 15 & [-5.000, 9.000] & 0.0e+00 \\
\bottomrule
\end{tabular}

\caption{Conteo de planos ocupados.}
\label{tab:espectral}
\end{table}

RANDU toma exactamente \textbf{15} valores, $k \in \{-5,\dots,9\}$, con
desviación $0.00$ respecto al entero más cercano: la relación es algebraica,
no aproximada. El LCG de Numerical Recipes toma 0 valores enteros y su $k_i$ se
distribuye de forma continua.

El conteo coincide con la predicción del Teorema~\ref{teo:red}: como cada
$u\in[0,1)$, la combinación $9u_i-6u_{i+1}+u_{i+2}$ está acotada entre $-6$ y
$10$, así que solo caben los 15 enteros de $-5$ a $9$. Se predijeron 15 y se
observaron 15.

\paragraph{Relación con el período y la calidad.} RANDU tiene período
$2^{29} \approx 5.4\times10^8$: más de quinientos millones de valores antes de
repetirse. Y sus ternas ocupan 15 planos. Son propiedades
\textbf{independientes}: el período mide cuánto tarda la secuencia en repetirse
y la prueba espectral mide qué tan bien llena el espacio. Un generador puede
tardar una eternidad en repetirse y distribuir pésimo en tres dimensiones.

Peor aún, la cota de Marsaglia permitía hasta
$(3!\cdot 2^{31})^{1/3} \approx 2344$ planos, que habría sido inofensivo.
RANDU se queda en 15: unas 150 veces por debajo de su propio peor caso
permitido.

\section{Tabla comparativa}
\label{sec:comparativa}

\begin{table}[H]
\centering\footnotesize
\begin{tabular}{lrrlp{4cm}p{4cm}}
\toprule
Algoritmo & s / $10^6$ & Rel. & Período & Seguridad & Casos de uso \\
\midrule
LCG (Num. Recipes) & 0.15 & 1.0$\times$ & $2^{32}$ & Nula: la forma cerrada permite recuperar $(a,c,m)$ con pocas salidas. & Simulación sencilla, videojuegos, sistemas embebidos, \texttt{rand()} de C. \\
RANDU & 0.15 & 1.0$\times$ & $2^{29}$ & Nula, y además las ternas caen en 15 planos. & Ninguno. Solo interés histórico y didáctico. \\
Cuadrados medios & 0.29 & 1.9$\times$ & cola $326\,565$ + ciclo $2\,500$ & Nula. Degenera: cinco puntos fijos con 4 dígitos. & Ninguno. Obsoleto desde los años 50. \\
Mersenne Twister & 0.66 & 4.4$\times$ & $2^{19937}-1$ & Nula: con 624 salidas se reconstruye el estado. & Estándar de simulación. Motor de \texttt{random} de Python. \\
Blum Blum Shub & 1.47 & 9.8$\times$ & $\approx 2^{51}$ & \textbf{Criptográfica}: equivale a la residuosidad cuadrática módulo $M$. & Criptografía, donde la lentitud es aceptable. \\
\bottomrule
\end{tabular}

\caption{Comparación de los cinco generadores. Las velocidades están
cronometradas en esta máquina (mejor de tres corridas); los períodos son
derivados donde hay teoría y medidos donde no la hay.}
\label{tab:comparativa}
\end{table}

El orden de velocidad reproduce el costo algebraico de cada recurrencia. El LCG
y RANDU son idénticos en costo, como debe ser: RANDU \emph{es} un LCG. Los
cuadrados medios pagan la conversión a cadena de dígitos. MT19937 paga la
torcedura de 624 palabras. Y BBS es diez veces más lento que el LCG porque
necesita varias exponenciaciones modulares para producir un solo número, ya que
solo puede emitir 5 bits por estado.

\subsection{Comparación con \texttt{random} y \texttt{secrets}}
\label{sec:librerias}

Para el inciso (d) se eligió MT19937 por una razón concreta: \texttt{random} de
Python \emph{es} MT19937. Al tratarse del mismo algoritmo, cualquier diferencia
mide la implementación y no la calidad del generador.

\begin{table}[H]
\centering\small
\begin{tabular}{lrcrrrr}
\toprule
Fuente & s / $10^6$ & Reproducible & $\chi^2$ & K--S & Autocorr. & Rachas \\
\midrule
MT19937 (propio) & 0.61 & Sí & 0.1142 & 0.0853 & 0.9599 & 0.9840 \\
random (MT19937 de C) & 0.05 & Sí & 0.3267 & 0.8335 & 0.8447 & 0.8729 \\
secrets (SystemRandom) & 0.24 & \textbf{No} & 0.6728 & 0.2604 & 0.8060 & 0.7949 \\
\bottomrule
\end{tabular}

\caption{Comparación del MT19937 propio contra las librerías estándar.}
\label{tab:librerias}
\end{table}

La implementación propia es unas doce veces más lenta que \texttt{random}, y
eso no dice nada sobre el algoritmo: es Python interpretado contra C compilado.
El detalle revelador es que \texttt{secrets} resulta \emph{más rápido} que el
MT19937 propio pese a ser criptográficamente seguro. Tampoco es una ventaja
algorítmica: \texttt{secrets} delega en el generador del sistema operativo, que
también está en C. Lo caro no es la criptografía, es el bucle de Python.

La diferencia que sí es de diseño está en la reproducibilidad.
\texttt{secrets} no acepta semilla, y no por omisión sino a propósito: un
generador criptográfico predecible no sirve de nada. Eso lo vuelve inutilizable
para simulación reproducible, que es justamente lo que este trabajo exige.
